% appB_apriori.tex: the a priori bound at the hard edge, and Q in I + dC(L^2).

\section{The a priori bound at the hard edge}\label{app:apriori}

This appendix proves item (vii) of Proposition~\ref{prop:frame}: $Q_\pm$ and $(Q^{-1})_\pm$ belong to
$I+\dC(L^2(\Sigma))$, together with a pointwise bound on $Q$ near $0$. The comparison argument of
Section~\ref{sec:hardedge} applies \cite[Theorem~4.4]{CD} (\S\ref{sec:tool-cd}(b)) to $Q$, and this membership is its
hypothesis: the boundary values of $Q-I$ must be Cauchy integrals of $L^2$ functions. The a priori bound of
Proposition~\ref{prop:apriori} holds for every $\alpha>-1$; the restriction $|\alpha|<1$ enters item (vii) only through
Corollary~\ref{cor:L2}, where the bound has to be square integrable at $0$.

We never solve a Riemann--Hilbert problem for $Q$; $Q$ is built from the orthogonal polynomials, and we only have to
check that it is not too singular at the vertices of $\Sigma$. The only delicate vertex is the hard edge $0$, where
infinitely many nodes accumulate. There we bound the Cauchy sums in $Y$ directly (Lemma~\ref{lem:sumbound}), except in
thin sectors around the band, where the poles are removed by $R$ and a contour argument is needed
(Lemma~\ref{lem:lenssector}). Throughout this appendix $n\ge n_0$ is \emph{fixed}, and constants may depend on $n$.

\subsection{Three facts about Cauchy integrals}\label{sec:B-prelim}

The following facts hold for general contours and are used here and in Appendix~\ref{app:resolvent}. In this
subsection $\Gamma$ is a bounded, finite union of oriented $C^2$ arcs with $|\Gamma\cap U(z,\rho)|\le C_\Sigma\rho$ for
all $z\in\C$, $\rho>0$, and $\mathcal V$ is the finite set of endpoints and intersection points of the arcs. We say that a
function $\mathfrak F$ analytic in $\C\setminus\Gamma$ is \emph{continuous up to $\Gamma$} if near every point of
$\Gamma\setminus\mathcal V$ its restriction to each side of $\Gamma$ extends continuously to the arc; $\mathfrak F_\pm$ then
denote these extensions. We use only Cauchy's integral formula, Green's formula and the facts of \S\ref{sec:tool-cd}(a).

\begin{lemma}[Cauchy integrals on a union of arcs]\label{lem:B-cauchy}
\begin{enumerate}
\item[(a)] For $\mathfrak h\in L^2(\Gamma)$ and $z\notin\Gamma$,
  \[
    |\Cauchy\mathfrak h(z)|\le\pi^{-1}C_\Sigma^{1/2}\|\mathfrak h\|_{L^2}\,\mathrm{dist}(z,\Gamma)^{-1/2}.
  \]
\item[(b)] Let $\mathfrak k$ be $C^1$ (in arc length) on each arc and vanish near $\mathcal V$. Then $\Cauchy\mathfrak k$ is
  analytic near each point of $\mathcal V$ and continuous up to $\Gamma$, and $(\Cauchy\mathfrak k)_+-(\Cauchy\mathfrak
  k)_-=\mathfrak k$ on $\Gamma\setminus\mathcal V$. In particular $(\Cauchy\mathfrak k)_\pm=\Cauchy_\pm\mathfrak k$.
\item[(c)] Let $\mathfrak F$ be analytic in $\C\setminus\Gamma$ and continuous up to $\Gamma$, with $\mathfrak F_\pm\in L^2(\Gamma)$,
  $\mathfrak F(z)=O(1/z)$ as $z\to\infty$, and $\rho\sup\{|\mathfrak F(z)|:|z-z_0|=\rho,\ z\notin\Gamma\}\to0$ as $\rho\to0$ for every
  $z_0\in\mathcal V$. Then $\mathfrak F=\Cauchy(\mathfrak F_+-\mathfrak F_-)$ in $\C\setminus\Gamma$; hence $\mathfrak F_\pm=\Cauchy_\pm(\mathfrak F_+-\mathfrak F_-)$, i.e.\
  $\mathfrak F_\pm\in\dC(L^2(\Gamma))$.
\end{enumerate}
\end{lemma}

\begin{proof}
(a) By Cauchy--Schwarz it suffices to show $\int_\Gamma|s-z|^{-2}|\diff s|\le4C_\Sigma/\mathfrak d$ with $\mathfrak d=
\mathrm{dist}(z,\Gamma)$. Split $\Gamma$ into the shells $\{2^k\mathfrak d\le|s-z|<2^{k+1}\mathfrak d\}$, $k\ge0$; the
$k$-th contributes at most $C_\Sigma2^{k+1}\mathfrak d\,(2^k\mathfrak d)^{-2}$, and the sum is $4C_\Sigma/\mathfrak d$.

(b) Analyticity near $\mathcal V$ is clear, since $\mathfrak k$ vanishes there. By a partition of unity we may assume
that $\mathfrak k$ is supported in a compact sub-arc of the interior of one arc $\gamma$. Choose an open disc $D$ which
meets $\Gamma$ only in an open sub-arc of $\gamma$ containing that support and which $\gamma$ cuts into two parts, $D_+$
on the left and $D_-$ on the right. Using normal coordinates along $\gamma$ (which are $C^1$ because $\gamma$ is $C^2$),
extend $\mathfrak k$ to a $C^1$ function $K$ on $\C$ with compact support in $D$ and $K|_\gamma=\mathfrak k$. Green's
formula on $D_+$ (the Cauchy--Pompeiu formula) gives, since the positively oriented boundary of $D_+$ runs along
$\gamma$ in its own direction and $K=0$ on $\partial D$,
\[
  \Cauchy\mathfrak k(z)=K(z)\mathbf 1_{D_+}(z)+\mathfrak P(z),\qquad
  \mathfrak P(z):=\frac1\pi\iint_{D_+}\frac{\bar\partial K(w)}{w-z}\diff A(w),\qquad z\notin\Gamma .
\]
Since $\bar\partial K$ is bounded with compact support and $|w-z|^{-1}$ is locally integrable uniformly in $z$, $\mathfrak P$
is continuous on $\C$. So $\Cauchy\mathfrak k$ is continuous up to $\gamma$ from both sides, with
$(\Cauchy\mathfrak k)_+-(\Cauchy\mathfrak k)_-=K=\mathfrak k$. As continuous extensions, these are the non-tangential
limits $\Cauchy_\pm\mathfrak k$ of \S\ref{sec:tool-cd}.

(c) Fix $z\notin\Gamma$. For small $\rho$ and large $R$ apply Cauchy's integral formula to $\mathfrak F$ on the open set
$U(0,R)\setminus\bigl(\Gamma\cup\bigcup_{z_0\in\mathcal V}\overline{U(z_0,\rho)}\bigr)$, whose boundary consists of the two sides
of the arcs of $\Gamma$ outside the small discs, the circles $|s-z_0|=\rho$ and the circle $|s|=R$; $\mathfrak F$ is continuous on
the closure (taken separately on the two sides of each arc). The two sides of an arc contribute
$\frac1{2\pi i}\int(\mathfrak F_+-\mathfrak F_-)(s)\diff s/(s-z)$, the left side being traversed in the direction of the arc. As $\rho\to0$
the small circles contribute at most $\sum_{z_0}2\rho\sup_{|s-z_0|=\rho}|\mathfrak F(s)|/|z-z_0|\to0$ (for $\rho\le|z-z_0|/2$), and the
arc integrals converge to $\Cauchy(\mathfrak F_+-\mathfrak F_-)(z)$ because $\mathfrak F_+-\mathfrak F_-\in L^2(\Gamma)\subset L^1(\Gamma)$. As $R\to\infty$ the large
circle contributes $O(1/R)$. Hence $\mathfrak F(z)=\Cauchy(\mathfrak F_+-\mathfrak F_-)(z)$. Taking boundary values from either side,
$\mathfrak F_\pm=\Cauchy_\pm(\mathfrak F_+-\mathfrak F_-)$ wherever the non-tangential limits $\Cauchy_\pm$ exist, i.e.\ almost everywhere.
\end{proof}

\subsection{The bound at the hard edge}\label{sec:B-bound}

Recall from Lemma~\ref{lem:interp} that
\begin{equation}\label{eq:U-sum}
  U_{j2}(z)=\sum_{k\ge1}\frac{w(X_k)\,U_{j1}(X_k)}{z-X_k}\qquad(j=1,2),\qquad w(X_k)=n^pw_k,
\end{equation}
and that $U_{11},U_{21}$ are polynomials. Put $\varphi_1:=\frac12\thlip$.

\begin{lemma}[sum bound]\label{lem:sumbound}
For every $\eta>0$ there is $C$ (depending on $n$ and $\eta$) such that $|U_{j2}(z)|\le C|z|^{-\eta}(1+|z|^{\alpha})$ for
$0<|z|\le1$, $\varphi_1\le|\arg z|\le\pi$.
\end{lemma}

\begin{proof}
For such $z$ and $X>0$, $|z-X|\ge\sin\varphi_1\max(|z|,X)$. By Lemma~\ref{lem:pottersector}(i),
$w_k\le C_\eta k^{-s+\eta}$, and $U_{j1}$ is bounded on $[0,n^p]$. Let $t_z:=|t(z)|=n|z|^{-1/p}$, so that $X_k\ge|z|$ if and only if
$k\le t_z$. Then
\[
  |U_{j2}(z)|\le C\sum_{k\le t_z}\frac{k^{-s+\eta}}{(n/k)^p}+\frac C{|z|}\sum_{k>t_z}k^{-s+\eta}
  \le C\bigl(1+t_z^{1+p-s+\eta}+\log t_z\bigr)+C|z|^{-1}t_z^{1-s+\eta}.
\]
Since $t_z^{1+p-s}=n^{1+p-s}|z|^{\alpha}$ and $|z|^{-1}t_z^{1-s}=n^{1-s}|z|^{\alpha}$ (recall $(s-1)/p-1=\alpha$), and
$t_z^\eta=n^\eta|z|^{-\eta/p}$, the claim follows after renaming $\eta$.
\end{proof}

\begin{lemma}[lens sectors]\label{lem:lenssector}
For every $\eta>0$ there are $\delta_{\mathrm s}\in(0,\delta)$ and $C$ such that $|R_{j2}(z)|\le C|z|^{-\eta}(1+|z|^\alpha)$ for
$0<|z|\le\delta_{\mathrm s}/2$, $0<|\arg z|<\varphi_1$, and also for $z\in\Ob\cap U(0,\delta_{\mathrm s}/2)$.
\end{lemma}

\begin{proof}
Inside the lens the Cauchy sum~\eqref{eq:U-sum} has poles that $R$ removes; we move the sum onto the boundary of a
sector by the residue theorem and bound the boundary integral. Take $z$ with $0<\arg z<\varphi_1$; the case $\arg z<0$ is the mirror image. In $\Ob^+$, $R_{j2}=U_{j2}-U_{j1}We^{-2i\theta}\upsilon_+$
by~\eqref{eq:R}. Put
\[
  \mathcal K_j(s):=\frac{U_{j1}(s)W(s)}{e^{2i\theta(s)}-1},\qquad s\in\mathcal G:=\{s:\ r<|s|<\delta_{\mathrm s},\ |\arg s|<\thlip\},
\]
where $\delta_{\mathrm s}$ and $r$ are chosen with $t(\delta_{\mathrm s}),t(r)\in\N+\frac12$, and $r<\frac12|z|$, $|z|<\frac12\delta_{\mathrm s}$. Since
$e^{-2i\theta}\upsilon_+=1/(e^{2i\theta}-1)$, we have $R_{j2}(z)=U_{j2}(z)-\mathcal K_j(z)$. The function $\mathcal K_j$ is meromorphic on a neighbourhood of $\overline{\mathcal G}$: $W$ is
analytic there (Lemma~\ref{lem:weightdomain}), and $e^{2i\theta}=1$ only where $t\in\mathbb Z$, i.e.\ at the nodes (the points with
$t\in-\N$ have $|\arg s|=\pi p>\thlip$). At a node, $e^{2i\theta(s)}-1=2i\theta'(X_k)(s-X_k)(1+O(s-X_k))$, so
$\Res_{X_k}\mathcal K_j=U_{j1}(X_k)W(X_k)/(2i\theta'(X_k))=w(X_k)U_{j1}(X_k)$. The residue theorem on $\partial\mathcal G$ (positively
oriented) and~\eqref{eq:U-sum} give
\begin{equation}\label{eq:R-contour}
  R_{j2}(z)=\sum_{X_k\notin\mathcal G}\frac{w(X_k)U_{j1}(X_k)}{z-X_k}-\frac1{2\pi i}\oint_{\partial\mathcal G}\frac{\mathcal K_j(s)\diff s}{s-z}.
\end{equation}
We bound the terms. (1) The nodes with $X_k\ge\delta_{\mathrm s}$ are finitely many and at distance $\ge\delta_{\mathrm s}/2$ from $z$; those with $X_k\le r$
contribute at most $\frac2{|z|}\sum_{X_k\le r}w(X_k)|U_{j1}(X_k)|$, which tends to $0$ as $r\to0$. (2) On the ray
$\arg s=\thlip$, $|e^{2i\theta}-1|^{-1}=|e^{-2i\theta}|/|1-e^{-2i\theta}|\le2\exp(-2\pi n|s|^{-1/p}\sin(\thlip/p))$ by~\eqref{eq:imtheta},
which beats $|W(s)|\le C|s|^{\alpha-\eta}$ (Lemmas~\ref{lem:bandweight}, \ref{lem:pottersector}(ii)); and $|s-z|\ge|s|\sin(\thlip-\varphi_1)$.
So this part is bounded. (3) On the ray $\arg s=-\thlip$ we write $1/(e^{2i\theta}-1)=-1-e^{2i\theta}/(1-e^{2i\theta})$, where now
$|e^{2i\theta}|<1$ is super-small, so $\mathcal K_j=-U_{j1}W+(\text{bounded})$. Since $|s-z|\ge(|s|+|z|)/C$,
\[
  \int_0^{\delta_{\mathrm s}}\frac{|W(\rho e^{-i\thlip})|\diff\rho}{\rho+|z|}\le C\int_0^{\delta_{\mathrm s}}\frac{\rho^{\alpha-\eta}\diff\rho}{\rho+|z|}\le C|z|^{-\eta}(1+|z|^{\alpha}),
\]
because $\alpha-\eta>-1$. (4) On the arcs $|s|=r$ and $|s|=\delta_{\mathrm s}$, $|t|\in\N+\frac12$. For such $t=u+iv$ we have
$|e^{2\pi it}-1|\ge1/C>0$: if $|v|\ge\frac14$ this is clear, and if $|v|<\frac14$ then $u$ is at distance at least $\frac38$ from
$\mathbb Z$, so $|e^{2\pi it}-1|=2e^{-\pi v}|\sin\pi t|\ge2e^{-\pi/4}\sin(3\pi/8)$. Hence $|\mathcal K_j|\le C|W|\le C|s|^{\alpha-\eta}$
there. The arc $|s|=r$ has length $\le2\thlip r$ and $|s-z|\ge\frac12|z|$, so its contribution is $\le Cr^{1+\alpha-\eta}/|z|\to0$
as $r\to0$; the arc $|s|=\delta_{\mathrm s}$ contributes a bounded amount. Letting $r\to0$ in~\eqref{eq:R-contour} proves the bound for
$0<\arg z<\varphi_1$. For $z\in\Ob\cap U(0,\delta_{\mathrm s}/2)$ with $\varphi_1\le|\arg z|<\frac54\thlip$, $R_{j2}=U_{j2}-U_{j1}We^{\mp2i\theta}\upsilon_\pm$
is bounded by Lemma~\ref{lem:sumbound} plus a super-small term.
\end{proof}

\begin{proposition}[a priori bound]\label{prop:apriori}
Fix $n\ge n_0$. Then $\det Q\equiv1$, and:
\begin{enumerate}
\item[(i)] for every $\eta>0$ there is $C_{n,\eta}$ with $|Q(z)|\le C_{n,\eta}|z|^{-|\alpha|/2-\eta}$ for $z\in U(0,\delta)\setminus\Sigma$;
\item[(ii)] $Q$ is bounded on every bounded set at positive distance from $0$, and $Q(z)\to I$ as $z\to\infty$;
\item[(iii)] $Q$ has continuous boundary values on $\Sigma$ minus its vertices.
\end{enumerate}
The same holds for $Q^{-1}$.
\end{proposition}

\begin{proof}
We follow the transformations from $Y$ to $Q$ factor by factor; every factor except the Cauchy sums in $U$ is explicit,
and those sums are controlled by Lemmas~\ref{lem:sumbound} and~\ref{lem:lenssector}. First, $\det U=\det Y=1$
(Lemma~\ref{lem:interp}), and every other factor in~\eqref{eq:VM} and in $Q=D_\infty^{-\sigma_3}VD_W^{\sigma_3}$ has
determinant $1$; so $\det Q=1$ and $Q^{-1}$ has the same entries up to sign and order.

(i) Let $\eta>0$ and $|z|<\delta_{\mathrm s}/2$ (the annulus $\delta_{\mathrm s}/2\le|z|<\delta$ is covered by~(ii)). By~\eqref{eq:VM} and~\eqref{eq:Q},
$Q$ is a constant diagonal matrix times $U\,M\,D_W^{\sigma_3}$. We use: $|e^{\pm n\varphi}|\le e^{Cn|z|^{1/2}}\le C$ (Lemma~\ref{lem:signs}(iii)); $|F|^{\pm1}\le C$ outside
the lens and $|\tilde F|^{\pm1}\le C$ in the lens near $0$ (Definition~\ref{def:outer-admissible}(e)); $|D_W(z)|^{\pm1}\le C|z|^{\pm\alpha/2-\eta}$
(Lemma~\ref{lem:szego}(vi)); $|1/W(z)|\le C|z|^{-\alpha-\eta}$ in the lens (Lemmas~\ref{lem:bandweight},
\ref{lem:pottersector}(ii)); and $U_{j1}$ bounded.

\emph{Outside the lens} ($z\in\Oinf$), $|\arg z|\ge\frac34\thlip>\varphi_1$ and $M=e^{n\varphi\sigma_3}F^{\sigma_3}$. The first column of
$Q$ is a constant times $U_{\cdot1}e^{n\varphi}FD_W$, of size $\le C|z|^{\alpha/2-\eta}$; the second is a constant times
$U_{\cdot2}e^{-n\varphi}F^{-1}D_W^{-1}$, of size $\le C|z|^{-\eta}(1+|z|^\alpha)|z|^{-\alpha/2-\eta}$ by Lemma~\ref{lem:sumbound}. Both are
$\le C|z|^{-|\alpha|/2-2\eta}$.

\emph{In the lens} $\Ob^\pm$, $UM=R\,e^{n\varphi\sigma_3}A_\pm\tilde F^{\sigma_3}$, and
$e^{n\varphi\sigma_3}A_\pm=\bigl(\begin{smallmatrix}e^{n\varphi}&0\\\mp e^{n\varphi}/W&e^{-n\varphi}\end{smallmatrix}\bigr)$. The first column of
$R$ is that of $U$, and the second is bounded by Lemma~\ref{lem:lenssector}. Hence the first column of $Q$ is at most
$C(1+|z|^{-\alpha-2\eta}(1+|z|^\alpha))|z|^{\alpha/2-\eta}\le C|z|^{-|\alpha|/2-3\eta}$, and the second is at most
$C|z|^{-\eta}(1+|z|^\alpha)|z|^{-\alpha/2-\eta}\le C|z|^{-|\alpha|/2-2\eta}$. Renaming $\eta$ gives~(i).

(ii) Away from $0$ each of $U$ (meromorphic, with finitely many nodes in any compact set not containing $0$), the triangular
and diagonal factors in~\eqref{eq:R}--\eqref{eq:V}, $F$, $\tilde F$ and $D_W$ is analytic in each region up to its boundary, except
for the poles of $U$ and of $\upsilon_\pm$ at the nodes, which cancel in $R$ (Step~1 of \S\ref{sec:frame-transf}). Near $4c$ the
only singular factors are $\tilde F\sim(4c-z)^{-1/2}$ and $D_W\sim(4c-z)^{1/2}$, and they occur only in the product
$\tilde FD_W$, whose square $(\tilde FD_W)^2=\mathsf w_nz^\alpha r_n\omega$ is bounded above and below on $U(4c,\eps)$
(Lemma~\ref{lem:szego}(vi)(b)). At the remaining vertices ($4c\pm ih_0$, $n^p+\delta'\pm ih_0$, $n^p+\delta'$) all factors are
analytic and bounded. At $\infty$, $Q\to I$ by Lemma~\ref{lem:normalisation}. (iii) follows from the same description and
Lemma~\ref{lem:jumpsV}.
\end{proof}

\begin{corollary}[$Q$ is an admissible solution]\label{cor:L2}
For each fixed $n\ge n_0$, $Q_\pm-I\in L^2(\Sigma)$ and $Q=I+\Cauchy(Q_+-Q_-)$ off $\Sigma$. Thus $Q_\pm\in I+\dC(L^2(\Sigma))$ in
the sense of \S\ref{sec:tool-cd}, and the same holds for $Q^{-1}$.
\end{corollary}

\begin{proof}
We apply Lemma~\ref{lem:B-cauchy}(c) to $\mathfrak F:=Q-I$ on $\Gamma=\Sigma$, which for fixed $n$ is bounded and satisfies
$|\Sigma\cap U(z,\rho)|\le C\rho$ (Proposition~\ref{prop:frame}(i)). By Proposition~\ref{prop:apriori}(iii), $\mathfrak F$ is continuous up to $\Sigma$. By
Proposition~\ref{prop:apriori}(i) with $\eta=(1-|\alpha|)/4$, which is positive because $|\alpha|<1$, we have
$|Q(z)|\le C|z|^{-(1+|\alpha|)/4}$ near $0$, and $(1+|\alpha|)/4<\frac12$; with~(ii) this gives $\mathfrak F_\pm\in L^2(\Sigma)$ and
$\rho\sup_{|z|=\rho}|\mathfrak F(z)|=O(\rho^{1-(1+|\alpha|)/4})\to0$ at the vertex $0$, while at every other vertex the
corresponding quantity is $O(\rho)$. Finally $\mathfrak F=O(1/z)$ at $\infty$ (Lemma~\ref{lem:normalisation}). The argument for
$Q^{-1}$ is identical.
\end{proof}

\noindent Proposition~\ref{prop:apriori} and Corollary~\ref{cor:L2} together are item (vii) of
Proposition~\ref{prop:frame}. Any $\eta\in(0,\frac12(1-|\alpha|))$ in Proposition~\ref{prop:apriori}(i) makes the
boundary values of $Q$ square integrable at $0$; this is the use (a) of $|\alpha|<1$ in Remark~\ref{rem:alpha-ge-1}.
